\section{Objectives and Endpoints}

This study pairs each objective with a prospectively defined, measurable endpoint
and a scientific justification, presented in the standard three-column form. Unlike
the single-arm first-in-human predicate \cite{kawchak2026phase1}, this protocol is
a randomized controlled efficacy study that directly compares Arm A (perioperative
daraxonrasib (RMC-6236) at the recommended Phase 2 dose (RP2D) plus on-premises
large language model (LLM) directed eight-arm robotic pancreaticoduodenectomy
(Whipple), PancreSpeed II) against Arm B (modified FOLFIRINOX
\cite{Conroy2018} plus institutional-standard high-volume pancreaticoduodenectomy).
The objective set is therefore organized around a single confirmatory efficacy
question, a fixed-sequence hierarchy of key secondary questions that share the
study-wide error budget, a set of estimation-only secondary measures, and a tier
of exploratory measures. Table~\ref{tab:objend} is the governing objectives-endpoints
table; it drives the analysis plan in \S9 and the safety rules in \S6 and \S7.

\noindent \textbf{Abbreviations in Trial Protocol:} BICR, blinded independent
central review; ctDNA, circulating tumor DNA; EBL, estimated blood loss; HR,
hazard ratio; ISGPS, International Study Group on Pancreatic Surgery; LLM, large
language model; LOS, length of stay; MPR, major pathologic response; OS, overall
survival; PDAC, pancreatic ductal adenocarcinoma; PFS, progression-free survival;
QoL, quality of life; RECIST, Response Evaluation Criteria in Solid Tumors; RP2D,
recommended Phase 2 dose. The internal control is Arm B; the 2025 Dutch nationwide
cohort of 1000 robotic pancreaticoduodenectomies \cite{DutchCohort2025} is retained
only as an external descriptive benchmark.

\vspace{0.3em}

\begin{table}[htbp]
\centering
\footnotesize
\caption{Objectives, endpoints, and justifications. The primary objective is
confirmatory; the key secondary objectives are tested in a fixed-sequence
hierarchy that controls the family-wise error rate; secondary objectives are
estimation only; exploratory objectives are hypothesis-generating.}
\label{tab:objend}
\begin{tabularx}{\textwidth}{L{4cm} Y Y}
\toprule
\textbf{Objective} & \textbf{Endpoint} & \textbf{Justification} \\
\midrule
\multicolumn{3}{l}{\textit{Primary}} \\
\midrule
Compare progression-free survival (PFS) between Arm A and Arm B. & PFS, defined as time from randomization to RECIST 1.1 progression or death from any cause, assessed by the investigator with blinded independent central review (BICR) as a sensitivity analysis. & PFS is the standard randomized Phase 2 efficacy anchor in resectable and borderline-resectable PDAC; a hazard ratio favoring Arm A supports progression to a Phase 3 confirmatory trial. \\
\midrule
\multicolumn{3}{l}{\textit{Key Secondary (hierarchical, fixed sequence)}} \\
\midrule
1. Compare overall survival (OS) between arms. & OS, time from randomization to death from any cause. & OS is the definitive measure of clinical benefit and the first gatekeeper after PFS in the testing hierarchy. \\
\addlinespace
2. Compare the R0 resection rate between arms. & Proportion with a microscopically margin-negative (R0) resection on central masked pathology review. & R0 status is the strongest modifiable surgical determinant of survival in resectable PDAC. \\
\addlinespace
3. Compare clinically relevant postoperative pancreatic fistula between arms. & ISGPS grade B or C postoperative pancreatic fistula rate \cite{Bassi2017}. & Grade B/C fistula is leading driver of post-Whipple morbidity + key safety quality discriminator vs.\ 2 operations. \\
\addlinespace
4. Compare major pathologic response between arms. & Major pathologic response (MPR) rate on central masked pathology review. & MPR is a validated surrogate of neoadjuvant treatment effect that links systemic regimen to resection specimen. \\
\addlinespace
5. Compare molecular response between arms. & KRAS circulating tumor DNA (ctDNA) clearance at week 12 \cite{Tie2019ctdna}. & Week-12 ctDNA clearance is an early, minimally invasive molecular readout of disease control. \\
\midrule
\multicolumn{3}{l}{\textit{Secondary (estimation)}} \\
\midrule
Estimate major postoperative complications by arm. & Clavien-Dindo grade III or higher complication rate through 90 days; external comparator 41.3\% \cite{DutchCohort2025}. & Clavien-Dindo grade III+ is the validated severity threshold for complications requiring procedural, endoscopic, or intensive-care management. \\
\addlinespace
Estimate early postoperative mortality by arm. & 90-day all-cause mortality. & Ninety-day mortality is the standard short-term survival safety measure for major hepatopancreatobiliary surgery. \\
\addlinespace
Estimate the conversion pathway by arm. & Unplanned conversion rate from the robotic system to an open or conventional approach; external comparator 10.1\% \cite{DutchCohort2025}. & Conversion rate contextualizes device performance against established robotic practice. \\
\addlinespace
Estimate operative burden and recovery by arm. & Operative time, estimated blood loss (EBL), and length of stay (LOS). & These operative + recovery measures characterize footprints of each operation. \\
\addlinespace
Estimate quality of life by arm. & EORTC QLQ-C30 and PAN26 scores over follow-up \cite{EORTCqlqc30}; 12- and 18-month PFS rates and 24-month OS rate. & Patient-reported quality of life and landmark survival rates frame the patient-centered value of each regimen. \\
\bottomrule \vspace{-0.2cm}
Table 4 continued, next page.\\
\end{tabularx}
\end{table}









\begin{table}[htbp]
\centering
\footnotesize
\begin{tabularx}{\textwidth}{L{4cm} Y Y}
\toprule
\textbf{Objective} & \textbf{Endpoint} & \textbf{Justification} \\
\midrule
\multicolumn{3}{l}{\textit{Exploratory}} \\
\midrule
Characterize the LLM advisory layer. & Concordance between the LLM advisory output and the independent hard-coded safety gate, as a proportion of advisories. & Advisory-to-gate concordance characterizes second-opinion oracle; informs later autonomy-graduation. \\
\addlinespace
Characterize simulation fidelity and the multicenter fleet. & Sim-to-real gap in trajectory (mm) and tip force (N); federated-learning model performance and ctDNA dynamics across sites. & These measures validate the upgraded Phase 0 gate and the cross-site fleet, supporting the digital-evidence basis for the device claim. \\
\addlinespace
Characterize health-economic value. & Cost and cost-effectiveness outcomes linked to the co-investment design. & Health-economic outcomes inform the equitable-access rationale and the capital-firewalled co-investment structure. \\
\bottomrule
\end{tabularx}
\end{table}














\vspace{0.3em}

Figure 7 renders the objective-to-endpoint hierarchy that organizes
Table~\ref{tab:objend}, showing how the single primary efficacy comparison of PFS
between Arm A and Arm B sits above the fixed-sequence key secondary family, which
in turn sits above the estimation-only secondary measures, exploratory
measures.

\begin{mermaidfig}
\node[mmgoal] (prim) at (0,-0.45)       {PRIMARY\\PFS: Arm A vs Arm B};
\node[mmstep] (sec)  at (0,-2.6)    {KEY SECONDARY\\(fixed sequence)\\OS; R0; ISGPS B/C;\\MPR; ctDNA};
\node[mmin]   (oth)  at (-4,-4.7) {SECONDARY\\(estimation)\\Clavien III+; mortality;\\conversion; QoL};
\node[mmin]   (exp)  at (4,-4.7)  {EXPLORATORY\\LLM concordance;\\sim-to-real; federated\\model; cost};
\draw[mmedged] (prim) -- (sec);
\draw[mmedged] (sec) -- (oth);
\draw[mmedged] (sec) -- (exp);
\end{mermaidfig}
\vspace{-0.7cm}
\figcaption{Figure 7. Objective-to-endpoint hierarchy for the randomized
comparison. The primary efficacy endpoint (progression-free survival, Arm A versus
Arm B) sits above the fixed-sequence key secondary family (overall survival, R0
resection rate, ISGPS grade B/C fistula, major pathologic response, and week-12
KRAS ctDNA clearance), which sits above the estimation-only secondary measures
(Clavien-Dindo grade III+, 90-day\\mortality, conversion rate, and quality of life)
and the exploratory measures (LLM advisory concordance,\\the sim-to-real gap,
federated model performance, and cost-effectiveness). Dashed edges separate the
tiers.}

\subsection{Primary Objective and Endpoint}

The single primary objective is to compare progression-free survival between Arm A
and Arm B. The primary endpoint is PFS, defined as the time from randomization to
the first documented RECIST 1.1 disease progression or death from any cause,
whichever occurs first. The primary analysis is investigator-assessed, with BICR
applied as a prespecified sensitivity analysis to guard against evaluation bias in
an open-label design. PFS is the established randomized Phase 2 efficacy anchor in
resectable and borderline-resectable PDAC: it is measurable within the trial
horizon, it integrates the systemic and surgical effects of each arm, and a hazard
ratio favoring Arm A provides the confirmatory signal needed to justify a Phase 3
trial. The study is sized for approximately 140 PFS events to provide 85\% power
at a two-sided significance level of 0.05 to detect a hazard ratio of 0.60
(median PFS 14.0 months in Arm B versus 23.3 months in Arm A), with a single
interim analysis for efficacy and futility at approximately 60\% of events
governed by an O'Brien-Fleming spending function \cite{DeMets1994lan,Jennison2000group}
(\S9).

\subsection{Key Secondary Objectives and Endpoints}

The key secondary objectives are tested in a single fixed-sequence (hierarchical)
gatekeeping procedure that controls the family-wise type I error rate at a
two-sided 0.05 level. Each hypothesis is tested at the full 0.05 level only if the
primary PFS comparison and every preceding key secondary hypothesis have reached
statistical significance; the first non-significant result closes the sequence, and
all subsequent comparisons are reported as descriptive only. The fixed order is:
(1) overall survival; (2) R0 resection rate; (3) ISGPS grade B/C postoperative
pancreatic fistula rate \cite{Bassi2017}; (4) major pathologic response rate; and
(5) week-12 KRAS ctDNA clearance \cite{Tie2019ctdna}. Overall survival is placed
first as the definitive measure of clinical benefit. R0 resection rate and MPR are
assessed on central masked pathology review, and the ISGPS fistula grade is
adjudicated centrally, so that each surgical-quality comparison is protected from
unmasking in the open-label setting. Week-12 ctDNA clearance provides an early
molecular corroboration of the imaging-based primary endpoint. This single ordered
sequence keeps the confirmatory inference for the entire key secondary family
within one shared error budget.

\subsection{Secondary Objectives and Endpoints}

The secondary objectives are estimation only; each is summarized with a point
estimate and a two-sided 95\% confidence interval (for between-arm contrasts, a
difference or hazard ratio with its interval), and none is part of the confirmatory
hierarchy. They comprise the \textbf{Clavien-Dindo grade III or higher
complication rate} through 90 days, for which the 2025 Dutch nationwide robotic
cohort reported 41.3\% \cite{DutchCohort2025}; \textbf{90-day all-cause
mortality}; the \textbf{unplanned conversion rate} from the robotic system to an
open or conventional approach, with a Dutch external comparator of 10.1\%
\cite{DutchCohort2025}; \textbf{operative time, estimated blood loss, and length of
stay}; and \textbf{quality of life} measured by the EORTC QLQ-C30 and PAN26
instruments \cite{EORTCqlqc30}, together with the 12- and 18-month PFS rates and
the 24-month OS rate. Because these measures sit outside the testing hierarchy,
the Dutch values are interpreted as external descriptive benchmarks and the
between-arm contrasts are interpreted as estimates rather than as confirmatory
findings; the internal Arm B comparator remains the primary reference frame.

\subsection{Exploratory Objectives and Endpoints}

The exploratory objectives generate hypotheses for later-phase study and are
explicitly non-confirmatory; no exploratory analysis is used to make a go or no-go
decision on the intervention, and all are reported with that caveat. \textbf{LLM
advisory-to-gate concordance} is the proportion of LLM advisory outputs that agree
with the independent hard-coded safety gate, quantifying the reliability of the
second-opinion oracle and informing autonomy-graduation decisions reserved for
later phases (\S4). The \textbf{sim-to-real gap} is the difference between the
Phase 0 simulation prediction and the intra-operative record, reported as the
trajectory gap in millimeters and the tip-force gap in newtons against the upgraded
Phase 0 thresholds, and is reported alongside \textbf{federated-learning model
performance} and \textbf{ctDNA dynamics} pooled across the eight sites to
characterize the multicenter fleet. The \textbf{health-economic and
cost-effectiveness} outcomes are linked to the capital-firewalled co-investment
design and inform the equitable-access rationale. Because these measures are
descriptive, they are summarized with point estimates and dispersion only, without
confirmatory significance testing.
